\IfFileExists{stacks-project.cls}{%
\documentclass{stacks-project}
}{%
\documentclass{amsart}
}

% For dealing with references we use the comment environment
\usepackage{verbatim}
\newenvironment{reference}{\comment}{\endcomment}
%\newenvironment{reference}{}{}
\newenvironment{slogan}{\comment}{\endcomment}
\newenvironment{history}{\comment}{\endcomment}

% For commutative diagrams we use Xy-pic
\usepackage[all]{xy}

% We use 2cell for 2-commutative diagrams.
\xyoption{2cell}
\UseAllTwocells

% We use multicol for the list of chapters between chapters
\usepackage{multicol}

% This is generally recommended for better output
\usepackage{lmodern}
\usepackage[T1]{fontenc}

% For cross-file-references
\usepackage{xr-hyper}

% Package for hypertext links:
\usepackage{hyperref}

% For any local file, say "hello.tex" you want to link to please
% use \externaldocument[hello-]{hello}
\externaldocument[introduction-]{introduction}
\externaldocument[conventions-]{conventions}
\externaldocument[sets-]{sets}
\externaldocument[categories-]{categories}
\externaldocument[topology-]{topology}
\externaldocument[sheaves-]{sheaves}
\externaldocument[sites-]{sites}
\externaldocument[stacks-]{stacks}
\externaldocument[fields-]{fields}
\externaldocument[algebra-]{algebra}
\externaldocument[brauer-]{brauer}
\externaldocument[homology-]{homology}
\externaldocument[derived-]{derived}
\externaldocument[simplicial-]{simplicial}
\externaldocument[more-algebra-]{more-algebra}
\externaldocument[smoothing-]{smoothing}
\externaldocument[modules-]{modules}
\externaldocument[sites-modules-]{sites-modules}
\externaldocument[injectives-]{injectives}
\externaldocument[cohomology-]{cohomology}
\externaldocument[sites-cohomology-]{sites-cohomology}
\externaldocument[dga-]{dga}
\externaldocument[dpa-]{dpa}
\externaldocument[sdga-]{sdga}
\externaldocument[hypercovering-]{hypercovering}
\externaldocument[schemes-]{schemes}
\externaldocument[constructions-]{constructions}
\externaldocument[properties-]{properties}
\externaldocument[morphisms-]{morphisms}
\externaldocument[coherent-]{coherent}
\externaldocument[divisors-]{divisors}
\externaldocument[limits-]{limits}
\externaldocument[varieties-]{varieties}
\externaldocument[topologies-]{topologies}
\externaldocument[descent-]{descent}
\externaldocument[perfect-]{perfect}
\externaldocument[more-morphisms-]{more-morphisms}
\externaldocument[flat-]{flat}
\externaldocument[groupoids-]{groupoids}
\externaldocument[more-groupoids-]{more-groupoids}
\externaldocument[etale-]{etale}
\externaldocument[chow-]{chow}
\externaldocument[intersection-]{intersection}
\externaldocument[pic-]{pic}
\externaldocument[weil-]{weil}
\externaldocument[adequate-]{adequate}
\externaldocument[dualizing-]{dualizing}
\externaldocument[duality-]{duality}
\externaldocument[discriminant-]{discriminant}
\externaldocument[derham-]{derham}
\externaldocument[local-cohomology-]{local-cohomology}
\externaldocument[algebraization-]{algebraization}
\externaldocument[curves-]{curves}
\externaldocument[resolve-]{resolve}
\externaldocument[models-]{models}
\externaldocument[functors-]{functors}
\externaldocument[equiv-]{equiv}
\externaldocument[pione-]{pione}
\externaldocument[etale-cohomology-]{etale-cohomology}
\externaldocument[proetale-]{proetale}
\externaldocument[relative-cycles-]{relative-cycles}
\externaldocument[more-etale-]{more-etale}
\externaldocument[trace-]{trace}
\externaldocument[crystalline-]{crystalline}
\externaldocument[spaces-]{spaces}
\externaldocument[spaces-properties-]{spaces-properties}
\externaldocument[spaces-morphisms-]{spaces-morphisms}
\externaldocument[decent-spaces-]{decent-spaces}
\externaldocument[spaces-cohomology-]{spaces-cohomology}
\externaldocument[spaces-limits-]{spaces-limits}
\externaldocument[spaces-divisors-]{spaces-divisors}
\externaldocument[spaces-over-fields-]{spaces-over-fields}
\externaldocument[spaces-topologies-]{spaces-topologies}
\externaldocument[spaces-descent-]{spaces-descent}
\externaldocument[spaces-perfect-]{spaces-perfect}
\externaldocument[spaces-more-morphisms-]{spaces-more-morphisms}
\externaldocument[spaces-flat-]{spaces-flat}
\externaldocument[spaces-groupoids-]{spaces-groupoids}
\externaldocument[spaces-more-groupoids-]{spaces-more-groupoids}
\externaldocument[bootstrap-]{bootstrap}
\externaldocument[spaces-pushouts-]{spaces-pushouts}
\externaldocument[spaces-chow-]{spaces-chow}
\externaldocument[groupoids-quotients-]{groupoids-quotients}
\externaldocument[spaces-more-cohomology-]{spaces-more-cohomology}
\externaldocument[spaces-simplicial-]{spaces-simplicial}
\externaldocument[spaces-duality-]{spaces-duality}
\externaldocument[formal-spaces-]{formal-spaces}
\externaldocument[restricted-]{restricted}
\externaldocument[spaces-resolve-]{spaces-resolve}
\externaldocument[formal-defos-]{formal-defos}
\externaldocument[defos-]{defos}
\externaldocument[cotangent-]{cotangent}
\externaldocument[examples-defos-]{examples-defos}
\externaldocument[algebraic-]{algebraic}
\externaldocument[examples-stacks-]{examples-stacks}
\externaldocument[stacks-sheaves-]{stacks-sheaves}
\externaldocument[criteria-]{criteria}
\externaldocument[artin-]{artin}
\externaldocument[quot-]{quot}
\externaldocument[stacks-properties-]{stacks-properties}
\externaldocument[stacks-morphisms-]{stacks-morphisms}
\externaldocument[stacks-limits-]{stacks-limits}
\externaldocument[stacks-cohomology-]{stacks-cohomology}
\externaldocument[stacks-perfect-]{stacks-perfect}
\externaldocument[stacks-introduction-]{stacks-introduction}
\externaldocument[stacks-more-morphisms-]{stacks-more-morphisms}
\externaldocument[stacks-geometry-]{stacks-geometry}
\externaldocument[moduli-]{moduli}
\externaldocument[moduli-curves-]{moduli-curves}
\externaldocument[examples-]{examples}
\externaldocument[exercises-]{exercises}
\externaldocument[guide-]{guide}
\externaldocument[desirables-]{desirables}
\externaldocument[coding-]{coding}
\externaldocument[obsolete-]{obsolete}
\externaldocument[fdl-]{fdl}
\externaldocument[index-]{index}

% Theorem environments.
%
\theoremstyle{plain}
\newtheorem{theorem}[subsection]{Theorem}
\newtheorem{proposition}[subsection]{Proposition}
\newtheorem{lemma}[subsection]{Lemma}

\theoremstyle{definition}
\newtheorem{definition}[subsection]{Definition}
\newtheorem{example}[subsection]{Example}
\newtheorem{exercise}[subsection]{Exercise}
\newtheorem{situation}[subsection]{Situation}

\theoremstyle{remark}
\newtheorem{remark}[subsection]{Remark}
\newtheorem{remarks}[subsection]{Remarks}

\numberwithin{equation}{subsection}

% Macros
%
\def\lim{\mathop{\mathrm{lim}}\nolimits}
\def\colim{\mathop{\mathrm{colim}}\nolimits}
\def\Spec{\mathop{\mathrm{Spec}}}
\def\Hom{\mathop{\mathrm{Hom}}\nolimits}
\def\Ext{\mathop{\mathrm{Ext}}\nolimits}
\def\SheafHom{\mathop{\mathcal{H}\!\mathit{om}}\nolimits}
\def\SheafExt{\mathop{\mathcal{E}\!\mathit{xt}}\nolimits}
\def\Sch{\mathit{Sch}}
\def\Mor{\mathop{\mathrm{Mor}}\nolimits}
\def\Ob{\mathop{\mathrm{Ob}}\nolimits}
\def\Sh{\mathop{\mathit{Sh}}\nolimits}
\def\NL{\mathop{N\!L}\nolimits}
\def\CH{\mathop{\mathrm{CH}}\nolimits}
\def\proetale{{pro\text{-}\acute{e}tale}}
\def\etale{{\acute{e}tale}}
\def\QCoh{\mathit{QCoh}}
\def\Ker{\mathop{\mathrm{Ker}}}
\def\Im{\mathop{\mathrm{Im}}}
\def\Coker{\mathop{\mathrm{Coker}}}
\def\Coim{\mathop{\mathrm{Coim}}}

% Boxtimes
%
\DeclareMathSymbol{\boxtimes}{\mathbin}{AMSa}{"02}

%
% Macros for moduli stacks/spaces
%
\def\QCohstack{\mathcal{QC}\!\mathit{oh}}
\def\Cohstack{\mathcal{C}\!\mathit{oh}}
\def\Spacesstack{\mathcal{S}\!\mathit{paces}}
\def\Quotfunctor{\mathrm{Quot}}
\def\Hilbfunctor{\mathrm{Hilb}}
\def\Curvesstack{\mathcal{C}\!\mathit{urves}}
\def\Polarizedstack{\mathcal{P}\!\mathit{olarized}}
\def\Complexesstack{\mathcal{C}\!\mathit{omplexes}}
% \Pic is the operator that assigns to X its picard group, usage \Pic(X)
% \Picardstack_{X/B} denotes the Picard stack of X over B
% \Picardfunctor_{X/B} denotes the Picard functor of X over B
\def\Pic{\mathop{\mathrm{Pic}}\nolimits}
\def\Picardstack{\mathcal{P}\!\mathit{ic}}
\def\Picardfunctor{\mathrm{Pic}}
\def\Deformationcategory{\mathcal{D}\!\mathit{ef}}
\title{Every 1-morphism of fibred categories factors through a fully faithful one with a left adjoint}
\author{The Stacks Project Authors}
\date{stacks-project:categories.tex:local-snapshot-20261012; GNU FDL 1.2 or later}
\begin{document}
\maketitle
\noindent\textit{Editorial note.} Counted span is the author's own proof environment, source lines 6900-6976, following the statement at 6881-6898.\par\medskip
\noindent\textit{Editorial note.} The adjointness bijection is marked 'follows immediately from the definitions'; the morphism description it unfolds is given in-span at 6906-6908.\par\medskip
\noindent\textit{Editorial note.} The proof uses no lemma references; only the chapter definitions.\par\medskip
\noindent\textit{Allowed standard prerequisites (editorial).}\begin{itemize}
\item Fibred categories and strongly cartesian morphisms. Definitions of fibred category and strongly cartesian morphism in this chapter. The verification that X' is fibred and that u, v, w are 1-morphisms of fibred categories runs on these definitions; F being a 1-morphism of fibred categories makes F(h*x) -> F(x) strongly cartesian.
\item Choice of pullbacks for the structure functors. Made at proof lines 6911-6912 following the chapter convention. Fixes the strongly cartesian lifts h*x -> x and h*y -> y used for the comparison morphism c and the lift (a, c).
\end{itemize}\medskip
\begin{lemma}
\label{lemma-ameliorate-morphism-fibred-categories}
Let $\mathcal{C}$ be a category. Let $F : \mathcal{X} \to \mathcal{Y}$
be a $1$-morphism of fibred categories over $\mathcal{C}$.
There exist $1$-morphisms of fibred categories over $\mathcal{C}$
$$
\xymatrix{
\mathcal{X} \ar@<1ex>[r]^u &
\mathcal{X}' \ar[r]^v \ar@<1ex>[l]^w & \mathcal{Y}
}
$$
such that $F = v \circ u$ and such that
\begin{enumerate}
\item $u : \mathcal{X} \to \mathcal{X}'$ is fully faithful,
\item $w$ is left adjoint to $u$, and
\item $v : \mathcal{X}' \to \mathcal{Y}$ is a fibred category.
\end{enumerate}
\end{lemma}

\begin{proof}
Denote $p : \mathcal{X} \to \mathcal{C}$ and $q : \mathcal{Y} \to \mathcal{C}$
the structure functors. We construct $\mathcal{X}'$ explicitly as follows.
An object of $\mathcal{X}'$ is a quadruple $(U, x, y, f)$ where
$x \in \Ob(\mathcal{X}_U)$, $y \in \Ob(\mathcal{Y}_U)$
and $f : y \to F(x)$ is a morphism in $\mathcal{Y}_U$.
A morphism $(a, b) : (U, x, y, f) \to (U', x', y', f')$ is given
by $a : x \to x'$ and $b : y \to y'$ with $p(a) = q(b) : U \to U'$
and such that $f' \circ b = F(a) \circ f$.

\medskip\noindent
Let us make a choice of pullbacks for both $p$ and $q$ and let us
use the same notation to indicate them.
Let $(U, x, y, f)$ be an object and let $h : V \to U$ be a morphism.
Consider the morphism $c : (V, h^*x, h^*y, h^*f) \to (U, x, y, f)$
coming from the given strongly cartesian maps $h^*x \to x$ and $h^*y \to y$.
We claim $c$ is strongly cartesian in $\mathcal{X}'$ over $\mathcal{C}$.
Namely, suppose we are given an object $(W, x', y', f')$ of $\mathcal{X}'$,
a morphism $(a, b) : (W, x', y', f') \to (U, x, y, f)$ lying over
$W \to U$, and a factorization $W \to V \to U$ of $W \to U$ through $h$.
As $h^*x \to x$ and $h^*y \to y$ are strongly cartesian we obtain morphisms
$a' : x' \to h^*x$ and $b' : y' \to h^*y$ lying over the given morphism
$W \to V$. Consider the diagram
$$
\xymatrix{
y' \ar[d]_{f'} \ar[r] & h^*y \ar[r] \ar[d]_{h^*f} & y \ar[d]_f \\
F(x') \ar[r] & F(h^*x) \ar[r] & F(x)
}
$$
The outer rectangle and the right square commute.
Since $F$ is a $1$-morphism of fibred categories the morphism
$F(h^*x) \to F(x)$ is strongly cartesian.
Hence the left square commutes by the universal property
of strongly cartesian morphisms. This proves that $\mathcal{X}'$
is fibred over $\mathcal{C}$.

\medskip\noindent
The functor $u : \mathcal{X} \to \mathcal{X}'$ is given by
$x \mapsto (p(x), x, F(x), \text{id})$. This is fully faithful.
The functor $\mathcal{X}' \to \mathcal{Y}$ is given by
$(U, x, y, f) \mapsto y$. The functor $w : \mathcal{X}' \to \mathcal{X}$
is given by $(U, x, y, f) \mapsto x$. Each of these functors is
a $1$-morphism of fibred categories over $\mathcal{C}$ by our
description of strongly cartesian morphisms of $\mathcal{X}'$ over
$\mathcal{C}$. Adjointness of $w$ and $u$ means that
$$
\Mor_\mathcal{X}(x, x') =
\Mor_{\mathcal{X}'}((U, x, y, f), (p(x'), x', F(x'), \text{id})),
$$
which follows immediately from the definitions.

\medskip\noindent
Finally, we have to show that $\mathcal{X}' \to \mathcal{Y}$ is a fibred
category. Let $c : y' \to y$ be a morphism in $\mathcal{Y}$
and let $(U, x, y, f)$ be an object of $\mathcal{X}'$ lying over $y$.
Set $V = q(y')$ and let $h = q(c) : V \to U$. Let $a : h^*x \to x$
and $b : h^*y \to y$ be the strongly cartesian morphisms covering $h$.
Since $F$ is a $1$-morphism of fibred categories we may identify
$h^*F(x) = F(h^*x)$ with strongly cartesian morphism
$F(a) : F(h^*x) \to F(x)$. By the universal property
of $b : h^*y \to y$ there is a morphism $c' : y' \to h^*y$ in
$\mathcal{Y}_V$ such that $c = b \circ c'$. We claim that
$$
(a, c) : (V, h^*x, y', h^*f \circ c') \longrightarrow (U, x, y, f)
$$
is strongly cartesian in $\mathcal{X}'$ over $\mathcal{Y}$. To see this
let $(W, x_1, y_1, f_1)$ be an object of $\mathcal{X}'$, let
$(a_1, b_1) : (W, x_1, y_1, f_1) \to (U, x, y, f)$ be a morphism
and let $b_1 = c \circ b_1'$ for some morphism $b_1' : y_1 \to y'$.
Then
$$
(a_1', b_1') : (W, x_1, y_1, f_1) \longrightarrow (V, h^*x, y', h^*f \circ c')
$$
(where $a_1' : x_1 \to h^*x$ is the unique morphism lying over the
given morphism $q(b_1') : W \to V$ such that $a_1 = a \circ a_1'$)
is the desired morphism.
\end{proof}

\end{document}
